%% ======================================================================
\section{Method}
\label{sec:method}
%% ======================================================================
%% PREAMBLE REQUIREMENTS (move these four lines to main.tex):
%%   \usepackage{tikz}
%%   \usetikzlibrary{arrows.meta,positioning}
%%   \usepackage{pifont}
%%   \usepackage{xcolor}
%% and the three glyph macros below.
\definecolor{okgreen}{RGB}{0,122,87}
\definecolor{okred}{RGB}{176,42,32}
\newcommand{\yes}{\textcolor{okgreen}{\ding{51}}}
\newcommand{\no}{\textcolor{okred}{\ding{55}}}
\newcommand{\na}{\textcolor{black!45}{--}}

We describe the method in the order it runs: first how a perturbation pair is
built and frozen, then the four stages that turn four samples into one
committed answer, then the selectors we compare against.
Figure~\ref{fig:pipeline} gives the pipeline and its cost;
Table~\ref{tab:methods} summarises how the selectors differ.

\begin{figure*}[t]
\centering
\definecolor{cshared}{RGB}{232,232,232}
\definecolor{ccode}{RGB}{205,228,243}
\definecolor{ccall}{RGB}{252,230,192}
\definecolor{ccallb}{RGB}{191,118,0}
\definecolor{cout}{RGB}{206,238,227}
\resizebox{\textwidth}{!}{%
\begin{tikzpicture}[
  font=\small,
  base/.style={draw, rounded corners=2.5pt, align=center, inner sep=5pt,
               minimum height=12mm, minimum width=23mm},
  input/.style={base, fill=white},
  shared/.style={base, fill=cshared},
  code/.style={base, fill=ccode},
  call/.style={base, fill=ccall, draw=ccallb, line width=1pt},
  out/.style={base, fill=cout},
  ar/.style={-{Latex[length=2.2mm]}, semithick},
  lbl/.style={font=\scriptsize, align=center, inner sep=2.5pt},
]
\node[input]  (prob) at (0,0)    {Problem $x$};
\node[shared] (samp) at (3.5,0)  {$4$ CoT samples\\[-2pt]{\scriptsize $T{=}0.7$, seeds $0$--$3$}};
\node[code]   (gate) at (7.4,0)  {\textbf{1.} disagreement\\[-2pt]gate};
\node[call]   (ver)  at (11.3,0) {\textbf{2.} contrastive\\[-2pt]verification};
\node[code]   (rule) at (15.0,0) {\textbf{3.} strict\\[-2pt]decision rule};
\node[out]    (ans)  at (18.4,0) {committed\\[-2pt]answer};

\draw[ar] (prob) -- node[lbl, above] {$4$ solver\\[-3pt]calls} (samp);
\draw[ar] (samp) -- node[lbl, above] {$0$ calls} (gate);
\draw[ar] (gate) -- node[lbl, above] {$K \ge 2$\\[-3pt]$1$ checker call} (ver);
\draw[ar] (ver)  -- node[lbl, above] {$0$ calls} (rule);
\draw[ar] (rule) -- (ans);

%% unanimous bypass, routed orthogonally below the main line
\coordinate (b1) at ([yshift=-1.6cm]gate.south);
\coordinate (b2) at ([yshift=-1.6cm]ans.south);
\draw[ar, dashed] (gate.south) -- (b1)
  -- node[lbl, above] {$K = 1$ (unanimous): keep the majority answer, no checker call} (b2)
  -- (ans.south);
\end{tikzpicture}%
}
\caption{The ESV pipeline, shaded by what each step costs. Stage~0 (grey) draws
four chain-of-thought samples, and those samples are written once and read
unchanged by every selector in Table~\ref{tab:methods}. ESV adds three stages:
two are deterministic code and free (blue), and only contrastive verification
issues a model call (orange, heavy border) --- one call, and only on the
instances where the samples disagree. Unanimous instances take the dashed path
and cost nothing beyond sampling. Because ESV chooses among candidates the
solver already produced, CoT pass@4 is a hard ceiling on the committed answer.}
\label{fig:pipeline}
\end{figure*}

\subsection{Constructing the perturbation pairs}
\label{sec:construction}

Each pair is produced by one generator model and accepted only on the
agreement of two verifier models, independent of the generator and of each
other. None of the three is a solver in any later experiment. One of the two
verifiers is, however, the same model we later use as the independent checker,
so every perturbed problem in the dataset is one that checker answered
correctly during construction. We report the originals separately for this
reason: they carry no such acceptance filter, and the checker's gain over the
vote has to survive on them as well (Section~\ref{sec:results}).

\paragraph{Generation.}
Given an original problem $P$ and its reference answer, the generator names the
solution strategy $S$ that $P$ invites, alters exactly one condition, and
returns the perturbed problem $P'$ together with its new answer, the condition
it changed, the reason the change matters, an alternative strategy, the answer
$S$ would produce if applied unchanged, and a confidence level. A balancer
assigns each problem one mode --- strategy-breaking or strategy-preserving ---
before the first attempt and holds it across retries, which keeps the two modes
near parity rather than letting the generator drift toward whichever is easier
to produce.

\paragraph{Verification.}
Acceptance turns on two calls to each verifier, in an order chosen so that
anchoring is impossible.

\begin{description}
\item[\textnormal{\emph{Blind solve.}}] The verifier sees $P'$ alone --- not
  $P$, not $S$, and not the generator's answer --- and solves it. The pair is
  rejected if the verifier declares $P'$ ill-posed, if its answer disagrees
  with the generator's, or if its confidence falls below $0.6$. Withholding $P$
  is what makes the resulting answer an independent check rather than a
  ratification.
\item[\textnormal{\emph{Mechanical strategy application.}}] The verifier now
  sees $P$, $P'$, $S$ and the verified answer, and applies $S$ to $P'$ without
  adapting it. In breaking mode the pair is confirmed when the resulting answer
  differs from the verified answer; in preserving mode, when it agrees. A
  report that $S$ cannot be carried out at all counts as a soft confirmation of
  breaking.
\end{description}

A pair is frozen only if neither verifier objects and at least one returns a
concrete confirmation; a soft confirmation alone is never sufficient. Rejected
attempts are returned to the generator with the objection attached and retried
up to a cap, and every attempt is logged so that the frozen set can be rebuilt
from the log.

One accepted breaking pair illustrates the structure. The original asks how far
Sofie walks if she walks twice as far as Brian \emph{over seven days}, with
answer $35$; the perturbation changes the comparison to twice as far as Brian
walks \emph{in one day}, with answer $245$. The original strategy still
completes on $P'$ and still returns $35$, so $35$ is the trap, and any candidate
equal to it is wrong for a reason we recorded when the pair was built.

\subsection{Explicit Strategy Verification}
\label{sec:esv}

ESV adds three stages to a sampling step it shares with every other method, and
spends at most one model call on any instance.

\paragraph{Stage~0: shared sampling.}
Each solver answers each problem $n = 4$ times with chain-of-thought prompting
\cite{wei2022chain}, one call per sample at temperature $0.7$ under seeds $0$
to $3$, so that the four samples genuinely differ rather than repeating one
trajectory. Those four calls are the sampling budget against which every
later cost in this paper is measured. The samples are drawn once, written to
disk, and read unchanged by every selector, so no comparison below is affected
by sampling variance between methods.

\paragraph{Stage~1: the disagreement gate.}
\label{sec:gate}
The four answers are grouped into distinct candidates under a type-aware
normalisation, in order of first appearance, with unreadable answers dropped.
If fewer than two candidates survive, the samples already agree, verification
has nothing to arbitrate, and the majority answer is committed with no call.
The gate is free and gold-free: it reads the samples only.

\paragraph{Stage~2: contrastive verification.}
\label{sec:verify}
Every surviving candidate is verified in a single call, each represented by one
of its solutions. The prompt states that the solutions reach different answers,
so at most one can be correct and all may be wrong; it presents them in
arbitrary order; and it withholds the vote counts while saying explicitly that
multiplicity is not evidence. For each candidate the checker returns the
strategy in one sentence, then one line per given the strategy depends on ---
quoted word for word from the problem statement rather than from the solution's
restatement of it --- marked \textsc{matches} or \textsc{mismatch}, then one
line per condition the strategy requires, marked \textsc{holds} or
\textsc{violated}. It then locates the first step at which the solutions
diverge, and finally returns \textsc{sound} or \textsc{flawed} per candidate
with a failure type, the first flawed step quoted verbatim, and the fact that
step should have asserted. Verification runs at temperature $0$. Asking for
quoted evidence before a verdict is the design choice under test: it is meant
to force the checker to commit to specific, checkable claims about the problem
text instead of judging fluency.

\paragraph{Stage~3: the decision rule.}
\label{sec:decide}
Selection is deterministic code with no further model call. Write $u_k$ for
candidate $k$'s verdict, $\mu_k$ for its count of \textsc{mismatch} bindings
and $\nu_k$ for its count of \textsc{violated} conditions. The candidate is
\emph{strictly sound} only if all three agree:
\begin{equation}
  \mathcal{S} = \{\, k : u_k = \textsc{sound},\; \mu_k = 0,\; \nu_k = 0 \,\}.
  \label{eq:strict}
\end{equation}
If $\mathcal{S}$ holds exactly one candidate, that candidate is committed. If
it holds several, the most-voted among them wins, breaking ties toward the
majority answer and then toward first appearance. If it is empty, the majority
answer is kept. Requiring the verdict to be consistent with the checker's own
line-by-line findings costs nothing and discards verdicts the checker has
already contradicted.

\paragraph{Two consequences of this shape.}
ESV selects only among answers the solver already produced, so CoT pass@4 is a
hard ceiling and no new answer can be introduced. And because an empty
$\mathcal{S}$ returns the majority answer, ESV is a filter on self-consistency
rather than a replacement for it: it can only act where the samples already
disagree, and declines by default.

\paragraph{Checker conditions.}
We run stage~2 in two conditions that differ in nothing but which model is
called. Under \emph{self-check} the solver verifies its own samples; under
\emph{independent check} a single fixed checker verifies the samples of every
solver. Everything else --- prompt, candidates, order, temperature, decision
rule --- is held constant, so the contrast isolates the verifier.

\subsection{Methods compared}
\label{sec:comparisons}

Every selector consumes the stage~0 samples unchanged and differs only in the
evidence it reads off them. Table~\ref{tab:methods} sets out those differences.

\begin{table}[t]
\centering
\small
\setlength{\tabcolsep}{4pt}
\begin{tabular}{@{}lccccc@{}}
\toprule
 & DQA & maj@4 & USC & Judge & \textbf{ESV} \\
\midrule
Gold-free                      & \yes & \yes & \yes & \yes & \yes \\
Reuses the shared samples      & \no  & \yes & \yes & \yes & \yes \\
Can overturn the majority      & \na  & \no  & \yes & \yes & \yes \\
May decline to choose          & \na  & \no  & \no  & \yes & \yes \\
Votes withheld from the judge  & \na  & \na  & \no  & \yes & \yes \\
Quoted, structured evidence    & \no  & \no  & \no  & \no  & \yes \\
\midrule
Calls beyond sampling          & $4^{\dagger}$ & $0$ & $1$ & $1$ & $1$ \\
\bottomrule
\end{tabular}
\caption{What each method is allowed to see and do. Reading down the last two
columns shows the design of our control: judge-select matches ESV on every
property but one, so the comparison between them isolates ESV's structured form
from the checker's raw capability. CoT pass@4 is omitted because it is an oracle
over the candidate set rather than a method. $^{\dagger}$Direct QA draws its own
four samples from a no-reasoning prompt instead of reading the shared pool.}
\label{tab:methods}
\end{table}

\noindent\textbf{Direct QA} is the one arm that does not. It draws its own four
samples from a prompt that forbids reasoning and takes their majority,
establishing that the benchmarks are not solvable by recall.

\noindent\textbf{CoT pass@1} is the accuracy of a single sample, and
\textbf{CoT maj@4} \cite{wang2022self} the majority over four --- our baseline,
and the answer every selector must either confirm or overturn. \textbf{CoT
pass@4} marks an instance correct whenever any candidate is correct, giving the
oracle ceiling of Eq.~\eqref{eq:bounds}.

\noindent\textbf{Universal Self-Consistency} \cite{chen2023universal} is the
prior gold-free selector and our principal comparison. It receives all four
samples with duplicates shown, so consensus remains visible, and names the
sample it selects. Because an unresolved selection silently falls back to the
majority answer, we report the rate at which USC resolves a selection before
reporting its accuracy.

\noindent\textbf{Judge-select} is a form control, not a baseline. It is matched
to stage~2 on the checker, the candidates, their order, the withheld vote
counts, and the licence to reject every candidate, and it reasons freely before
naming its pick. Two of those choices deliberately favour the control: free-form
reasoning makes the comparison structured-versus-unstructured rather than
reasoning-versus-none, and the licence to decline denies ESV a free win on
instances where no candidate survives. Judge-select therefore isolates what
ESV's structured form contributes once the checker's raw capability is held
fixed.

Reference answers enter at scoring and nowhere else, for every method.
